\documentclass[11pt,spanish]{article}

\usepackage[utf8]{inputenc}
\usepackage{babel}
\usepackage{amsmath,amssymb}
\usepackage{xcolor}
\usepackage{graphicx}
\usepackage{upgreek}
\usepackage{fancyhdr}


%%%%%%%%%%%%%%%%%%%%%%%%%%%%%%%%%%%%%%%%%%%%%%
% Page layout by CBB 2023.03.13
%%%%%%%%%%%%%%%%%%%%%%%%%%%%%%%%%%%%%%%%%%%%%%

\oddsidemargin=0.0cm
\textwidth=17cm
\topmargin=-1.5cm
\textheight=51\baselineskip
\parindent 0pt

\pagestyle{fancy}
\fancyhf{}
\fancyhead[R]{\textcolor[rgb]{0.0,0.4,0.0}{MMFII 2026 -- Segundo Parcial}\hfill\textcolor[rgb]{0.0,0.4,0.0}{\thepage}}
\renewcommand{\headrulewidth}{.5pt}
\renewcommand{\headrule}{\hbox to\headwidth{\color[rgb]{0.0,0.4,0.0}{\leaders\hrule height \headrulewidth\hfill}}}

%%%%%%%%%%%%%%%%%%%%%%%%%%%%%%%%%%%%%%%%%%%%%%
% Added by CBB 2023.03.06
%%%%%%%%%%%%%%%%%%%%%%%%%%%%%%%%%%%%%%%%%%%%%%

\DeclareMathOperator{\real}{Re}
\DeclareMathOperator{\imag}{Im}
\DeclareMathOperator{\Span}{span}
\DeclareMathOperator{\rank}{rank}
\DeclareMathOperator{\diag}{diag}
\DeclareMathOperator{\rng}{rng}
\DeclareMathOperator{\trace}{Tr}
\DeclareMathOperator{\Arg}{arg}

\newcommand{\abs}[1]{\lvert{#1}\rvert}
\newcommand{\norm}[1]{\lVert{#1}\,\rVert}
\newcommand{\txt}[1]{\textnormal{#1}}
\newcommand{\cc}[1]{\bar{#1}}
\newcommand{\der}[2]{\frac{d {#1}}{d {#2}}}

\newcommand{\eps}{\varepsilon}

\newcommand{\sca}[1]{\mathbb{#1}}
\newcommand{\op}[1]{{\cal{#1}}}
\newcommand{\mat}[1]{\mathsf{#1}}
\newcommand{\gmat}[1]{\boldsymbol{#1}}
\newcommand{\cev}[1]{\reflectbox{$\vec{\reflectbox{$#1$}}$}}
\newcommand{\pvec}[1]{\vec{#1}{\,}'}
\newcommand{\svec}[1]{\vec{#1}{\,}\mbox{}}

\newcommand{\aplic}[3]{#1 : #2 \to #3}

\newcommand{\ud}[2]{^{#1}_{\phantom{#1}#2}}

\newcommand{\ul}[1]{{\underline{#1}}}

\newcommand{\inpr}[2]{\left\langle{#1},{#2}\right\rangle}
\newcommand{\matel}[3]{\left\langle{#1}\left\vert{#2}\right\vert{#3}\right\rangle}

%%%%%%%%%%%%%%%%%%%%%%%%%%%%%%%%%%%%%%%%%%%%%%
% Counters definitions (old)
%%%%%%%%%%%%%%%%%%%%%%%%%%%%%%%%%%%%%%%%%%%%%%

\newcounter{problema}

\newcommand{\prob}{\stepcounter{problema} \vspace{.5cm} \noindent{\bf Problema \arabic{problema}: }}
\renewcommand{\labelenumi}{\alph{enumi})}
\renewcommand{\labelenumii}{\arabic{enumii})}
\newcommand\al{\alpha}
\newcommand\noi{\noindent}

%%%%%%%%%%%%%%%%%%%%%%%%%%%%%%%%%%%%%%%%%%%%%%
%%%%%%%%%%%%%%%%%%%%%%%%%%%%%%%%%%%%%%%%%%%%%%

\begin{document}

%%%%%%%%%%%%%%%%%%%%%%%%%%%%%%%%%%%%%%%%%%%%%%
% Begin custom header - by CBB 2023.03.13
%%%%%%%%%%%%%%%%%%%%%%%%%%%%%%%%%%%%%%%%%%%%%%

\thispagestyle{empty} % no LaTeX-generated header for 1st page

\mbox{}\\
\begin{minipage}[c][0cm]{\textwidth}
%%%
\begin{tabular}{lcr}
%%

\includegraphics[scale=0.3,trim=0 0 40 0]{logo_unc.jpg}
%%
&
%%
\begin{minipage}[b][0cm]{0.5\textwidth}
%
\centering
{\large Métodos Matemáticos de la Física II} \\
\centering
{\large Examen Final -- 21/07/2026}  \\[-0.4\baselineskip]
\mbox{}
%
\end{minipage}
%%
&
%%

\includegraphics[scale=0.4,trim=80 65 0 0]{logo-FAMAF-2016.jpg}
%%
\end{tabular}
%%%
\\[5pt]
\hrule
\mbox{}\\[5\baselineskip]
%%%
\end{minipage}

\vspace{-0.7\baselineskip}
\begin{center}
\bf
Escriba su nombre en cada página y numere todas las páginas\\
Resuelva cada problema en hojas separadas\\
Se permite el uso de tablas de transformadas y de funciones especiales
\end{center}
%\vspace{-1.0\baselineskip}
%
\hrule


\vspace{5mm}

\prob{} La siguente matriz opera sobre un espacio vectorial $V$ de dimensión 4.
\[
A= \left( \begin{array}{rrrr}
-4 & 0 & -1 & 0 \\
0 & -2 &  0 & 0 \\
4 & 1 & 0 & 0 \\
6 & 1 & 3 & -2
\end{array} \right).
\]
Sabiendo que $det(A-\lambda I) = (\lambda+2)^4,$ Calcule:
\begin{enumerate}
\item Los autovectores de $A$.
\item La forma de Jordan de la matriz $A$ (no le pedimos el cambio de base)
\item Responda: >En cuántos subespacios invariantes irreducibles se descompone
$V$ bajo la acción de $A$? >Qué dimensión tiene cada uno de esos subespacios?
Justifique claramente.
\end{enumerate}

\vspace{0.5cm}
{\bf Soluci\'on }{\em 1)}
\vspace{0.5cm}

El único autovalor es
$$
\boxed{\lambda=-2}
$$
con multiplicidad algebraica $4$.

{\bf a) Autovectores}

Debemos resolver
$$
(A+2I)\mathbf v=0.
$$
Se tiene
$$
A+2I=
\begin{pmatrix}
-2&0&-1&0\\
0&0&0&0\\
4&1&2&0\\
6&1&3&0
\end{pmatrix}.
$$
Sea
$$
\mathbf v=(x,y,z,w)^T.
$$
Las ecuaciones son
$$
\begin{cases}
-2x-z=0,\\
4x+y+2z=0,\\
6x+y+3z=0.
\end{cases}
$$
De la primera,
$$
z=-2x.
$$
La segunda da
$$
4x+y-4x=0
\quad\Longrightarrow\quad
y=0.
$$
La tercera queda
$$
6x-6x=0,
$$
por lo que no aporta información nueva.
La coordenada $w$ es completamente libre.
Así,
$$
\mathbf v
=
x
\begin{pmatrix}
1\\0\\-2\\0
\end{pmatrix}
+
w
\begin{pmatrix}
0\\0\\0\\1
\end{pmatrix}.
$$
Una base del autoespacio es
$$
\boxed{
\left\{
\begin{pmatrix}
1\\0\\-2\\0
\end{pmatrix},
\begin{pmatrix}
0\\0\\0\\1
\end{pmatrix}
\right\}.
}
$$
Por lo tanto,
$$
\boxed{\dim E_{-2}=2.}
$$

{\bf b) Forma de Jordan}

La multiplicidad algebraica es
$$
m_a=4,
$$
mientras que la geométrica es
$$
m_g=2.
$$
Luego existen exactamente dos bloques de Jordan.
Las posibles particiones de $4$ en dos bloques son
\begin{itemize}
\item $3+1$,
\item $2+2$.
\end{itemize}
Para decidir cuál corresponde, estudiamos
$$
N=A+2I.
$$
Calculamos
$$
N^2=
\begin{pmatrix}
0&-1&0&0\\
0&0&0&0\\
0&2&0&0\\
0&4&0&0
\end{pmatrix},
$$
y
$$
N^3=0.
$$
Entonces
$$
N^2\neq0,\qquad N^3=0.
$$
El índice de nilpotencia es
$$
3,
$$
por lo que el mayor bloque de Jordan tiene tamaño $3$.
Como sólo existen dos bloques,
$$
3+1=4.
$$
Así, la forma de Jordan es
$$
\boxed{
J=
\begin{pmatrix}
-2&1&0&0\\
0&-2&1&0\\
0&0&-2&0\\
0&0&0&-2
\end{pmatrix}.
}
$$
Equivalentemente,
$$
J=
J_3(-2)\oplus J_1(-2).
$$

{\bf c) Subespacios invariantes irreducibles}

Cada bloque de Jordan corresponde a un subespacio invariante indecomponible.
Como la forma de Jordan posee
\begin{itemize}
\item un bloque de tamaño $3$,
\item un bloque de tamaño $1$,
\end{itemize}
el espacio vectorial se descompone como
$$
V=V_3\oplus V_1,
$$
donde
$$
\dim V_3=3,\qquad
\dim V_1=1.
$$
Por consiguiente,
$$
\boxed{\text{Existen exactamente dos subespacios invariantes irreducibles.}}
$$
Sus dimensiones son
$$
\boxed{3\quad\text{y}\quad1.}
$$
La justificación es que cada bloque de Jordan genera un subespacio cíclico invariante que no puede descomponerse en suma directa de subespacios invariantes más pequeños.

{\bf Respuestas finales}

\begin{itemize}
\item {\bf a)} Una base del autoespacio es
$$
\boxed{
\left\{
\begin{pmatrix}
1\\0\\-2\\0
\end{pmatrix},
\begin{pmatrix}
0\\0\\0\\1
\end{pmatrix}
\right\}.
}
$$

\item {\bf b)}
$$
\boxed{
J=J_3(-2)\oplus J_1(-2).
}
$$

\item {\bf c)}
$$
\boxed{
V=V_3\oplus V_1,\qquad
\dim V_3=3,\quad
\dim V_1=1.
}
$$
\end{itemize}









\prob{} La métrica sobre una superficie toroidal de radio grande $R$ y radio
pequeño $a$ ($a<R$), está dada en coordenadas $\{\theta, \alpha\}$ por
\[
ds^2 = (R + a \cos(\theta))^2 d\theta^2 + a^2 d\alpha^2,
\]
donde $\theta$ es el ángulo que mide la rotación alrededor al eje de simetría
axial del toroide y $\alpha$ mide la rotación sobre la circunferencia de radio
$a$ en un plano $\theta =$const. Calcule la métrica inversa $g^{ij}$, los símbolos de Christoffel y el
operador Laplaciano actuando sobre una función $u(\theta,\alpha)$, es decir
$\Delta u(\theta,\alpha) = g^{ij} \nabla_i \nabla_j u.$

\vspace{0.5cm}
{\bf Soluci\'on }{\em 2)}
\vspace{0.5cm}

La métrica del toro viene dada por
$$
ds^2=(R+a\cos\theta)^2\,d\theta^2+a^2\,d\alpha^2,
$$
donde $R>a>0$, y las coordenadas son
$$
x^1=\theta,\qquad x^2=\alpha.
$$
Debemos calcular:
\begin{itemize}
\item la métrica inversa $g^{ij}$,
\item los símbolos de Christoffel,
\item el Laplaciano
\end{itemize}
$$
\Delta u=g^{ij}\nabla_i\nabla_j u.
$$

{\bf 1. Tensor métrico}

La métrica es diagonal:
$$
g_{ij}=
\begin{pmatrix}
(R+a\cos\theta)^2 & 0\\
0 & a^2
\end{pmatrix}.
$$
Su determinante es
$$
g=\det(g_{ij})
=a^2(R+a\cos\theta)^2,
$$
de donde
$$
\sqrt g
=a(R+a\cos\theta).
$$

{\bf 2. Métrica inversa}

Al ser diagonal,
$$
g^{ij}=
\begin{pmatrix}
\dfrac1{(R+a\cos\theta)^2} & 0\\
0 & \dfrac1{a^2}
\end{pmatrix}.
$$
Es decir,
$$
\boxed{
g^{\theta\theta}
=\frac1{(R+a\cos\theta)^2},
\qquad
g^{\alpha\alpha}
=\frac1{a^2}.
}
$$

{\bf 3. Símbolos de Christoffel}

Recordemos
$$
\Gamma^k_{ij}
=
\frac12
g^{kl}
\left(
\partial_i g_{jl}
+\partial_j g_{il}
-\partial_l g_{ij}
\right).
$$
Como
$$
g_{\alpha\alpha}=a^2
$$
es constante,
$$
\partial_\theta g_{\alpha\alpha}=0,
\qquad
\partial_\alpha g_{ij}=0.
$$
Además,
$$
g_{\theta\theta}
=(R+a\cos\theta)^2,
$$
por lo que
$$
\partial_\theta g_{\theta\theta}
=
2(R+a\cos\theta)(-a\sin\theta).
$$
Es decir,
$$
\partial_\theta g_{\theta\theta}
=
-2a(R+a\cos\theta)\sin\theta.
$$

{\bf Cálculo de $\Gamma^\theta_{\theta\theta}$}

$$
\Gamma^\theta_{\theta\theta}
=
\frac12
g^{\theta\theta}
\partial_\theta g_{\theta\theta}.
$$
Entonces
$$
\frac12
\frac1{(R+a\cos\theta)^2}
(-2a(R+a\cos\theta)\sin\theta),
$$
de donde
$$
\boxed{
\Gamma^\theta_{\theta\theta}
=
-\frac{a\sin\theta}
{R+a\cos\theta}.
}
$$

{\bf Cálculo de $\Gamma^\theta_{\alpha\alpha}$}

Como
$$
g_{\alpha\alpha}
=\text{cte},
$$
resulta
$$
\boxed{
\Gamma^\theta_{\alpha\alpha}=0.
}
$$

{\bf Cálculo de $\Gamma^\alpha_{\theta\alpha}$}

$$
\Gamma^\alpha_{\theta\alpha}
=
\frac12
g^{\alpha\alpha}
\partial_\theta g_{\alpha\alpha}=0.
$$
Luego
$$
\boxed{
\Gamma^\alpha_{\theta\alpha}
=
\Gamma^\alpha_{\alpha\theta}
=0.
}
$$
Todos los demás símbolos también son nulos.
Por lo tanto,
$$
\boxed{
\Gamma^\theta_{\theta\theta}
=
-\frac{a\sin\theta}
{R+a\cos\theta},
}
$$
es el único símbolo distinto de cero.

{\bf 4. Laplaciano}

Para un escalar,
$$
\nabla_i u=\partial_i u,
$$
y
$$
\nabla_i\nabla_j u
=
\partial_i\partial_j u
-
\Gamma^k_{ij}\partial_k u.
$$
Así,
$$
\Delta u
=
g^{ij}\nabla_i\nabla_ju.
$$
Como sólo existe
$$
\Gamma^\theta_{\theta\theta},
$$
tenemos
$$
\Delta u
=
g^{\theta\theta}
\left(
u_{\theta\theta}
-
\Gamma^\theta_{\theta\theta}u_\theta
\right)
+
g^{\alpha\alpha}
u_{\alpha\alpha}.
$$
Sustituyendo,
$$
\frac1{(R+a\cos\theta)^2}
\left(
u_{\theta\theta}
+
\frac{a\sin\theta}
{R+a\cos\theta}
u_\theta
\right)
+
\frac1{a^2}
u_{\alpha\alpha}.
$$
Finalmente,
$$
\boxed{
\Delta u
=
\frac{1}{(R+a\cos\theta)^2}
\frac{\partial^2u}{\partial\theta^2}
+
\frac{a\sin\theta}
{(R+a\cos\theta)^3}
\frac{\partial u}{\partial\theta}
+
\frac1{a^2}
\frac{\partial^2u}{\partial\alpha^2}.
}
$$

{\bf Verificación mediante la expresión de Laplace–Beltrami}

Como control, usemos
$$
\Delta u
=
\frac1{\sqrt g}
\partial_i
\left(
\sqrt g\,g^{ij}\partial_j u
\right).
$$
Aquí,
$$
\sqrt g
=a(R+a\cos\theta),
$$
por lo que
$$
\sqrt g\,g^{\theta\theta}
=
\frac{a}{R+a\cos\theta},
$$
y
$$
\sqrt g\,g^{\alpha\alpha}
=
\frac{R+a\cos\theta}{a}.
$$
Entonces,
$$
\Delta u
=
\frac1{a(R+a\cos\theta)}
\left[
\partial_\theta
\left(
\frac{a}{R+a\cos\theta}
u_\theta
\right)
+
\partial_\alpha
\left(
\frac{R+a\cos\theta}{a}
u_\alpha
\right)
\right].
$$
Como el factor $(R+a\cos\theta)/a$ no depende de $\alpha$, se obtiene exactamente la misma expresión anterior, lo que confirma el resultado.

{\bf Resumen}

$$
\boxed{
g^{ij}
=
\begin{pmatrix}
\dfrac1{(R+a\cos\theta)^2} & 0\\
0 & \dfrac1{a^2}
\end{pmatrix},
}
$$
$$
\boxed{
\Gamma^\theta_{\theta\theta}
=
-\frac{a\sin\theta}{R+a\cos\theta},
}
$$
todos los demás símbolos de Christoffel son nulos, y
$$
\boxed{
\Delta u
=
\frac1{(R+a\cos\theta)^2}
u_{\theta\theta}
+
\frac{a\sin\theta}{(R+a\cos\theta)^3}
u_\theta
+
\frac1{a^2}
u_{\alpha\alpha}.
}
$$
Este resultado coincide con el operador de Laplace–Beltrami para la métrica dada.








\prob{} Resuelva mediante el método de las características
\[
2xy \frac{\partial u}{\partial x} + \frac{\partial u}{\partial y}  - u = 0,
\]
sabiendo que $u(x,y=0) = x.$

\vspace{0.5cm}
{\bf Soluci\'on }{\em 3)}
\vspace{0.5cm}

Se desea resolver la ecuación en derivadas parciales

\[
2xy\,\frac{\partial u}{\partial x}
+\frac{\partial u}{\partial y}
-u=0,
\]

con condición inicial

\[
u(x,0)=x.
\]

{\bf Método de las características}

Escribimos la ecuación como

\[
X(u)=u,
\]

donde

\[
X
=
2xy\,\frac{\partial}{\partial x}
+
\frac{\partial}{\partial y}
\]

es un campo vectorial sobre el plano $(x,y)$.

Introducimos ahora el espacio ampliado de coordenadas $(x,y,u)$ y consideramos el
campo vectorial

\[
\widetilde X
=
2xy\,\frac{\partial}{\partial x}
+
\frac{\partial}{\partial y}
+
u\,\frac{\partial}{\partial u}.
\]

Las curvas características son las curvas integrales de $\widetilde X$, es decir,
las soluciones del sistema

\[
\boxed{
\begin{aligned}
\dot x &=2xy,\\
\dot y &=1,\\
\dot u &=u,
\end{aligned}
}
\]

donde el punto denota derivación respecto del parámetro de la característica.

{\bf Cálculo de las características}

Como

\[
\dot y=1,
\]

podemos tomar

\[
y=s.
\]

La ecuación para $x$ queda

\[
\frac{dx}{dy}=2xy,
\]

que es una ecuación diferencial separable:

\[
\frac{dx}{x}=2y\,dy.
\]

Integrando,

\[
\ln|x|=y^2+C,
\]

o equivalentemente,

\[
\boxed{
xe^{-y^2}=C.
}
\]

Por lo tanto, cada característica queda identificada por el valor constante de

\[
\xi=xe^{-y^2}.
\]

{\bf Evolución del potencial}

Sobre cada característica,

\[
\frac{du}{dy}=u,
\]

cuya solución es

\[
u=A(\xi)e^y,
\]

donde $A$ depende únicamente de la constante característica
$\xi=xe^{-y^2}$.

En consecuencia,

\[
u(x,y)
=
e^yA\!\left(xe^{-y^2}\right).
\]

{\bf Condición inicial}

Imponiendo la condición

\[
u(x,0)=x,
\]

se obtiene

\[
A(x)=x.
\]

Por lo tanto,

\[
A(\xi)=\xi,
\]

y finalmente

\[
\boxed{
u(x,y)
=
x\,e^{\,y-y^2}.
}
\]

{\bf Verificación}

Calculamos

\[
u_x=e^{y-y^2},
\]

y

\[
u_y
=
x(1-2y)e^{y-y^2}.
\]

Entonces,

\[
2xy\,u_x+u_y-u
=
xe^{y-y^2}
\left[
2y+(1-2y)-1
\right]
=0,
\]

por lo que la ecuación diferencial queda satisfecha. Además,

\[
u(x,0)=x,
\]

verificando también la condición inicial.

\[
\boxed{
u(x,y)=x\,e^{\,y-y^2}.
}
\]








\prob{} El potencial $\Phi(r,\theta,z)$ satisface la ecuación de Laplace en el interior de un 
cilindro infinitamente largo de radio $b.$ En el borde del cilindro la función
satisface, para $V>0,$
\[
 \Phi(r=b, \theta,z) = \left\lbrace
\begin{array}{lcl}
 +V & \quad \mbox{si} \quad & 0<\theta\leq \pi/2 \\
-V & \mbox{si}   & \pi/2<\theta\leq \pi \\
+V & \quad \mbox{si} \quad & \pi<\theta\leq 3\pi/2 \\
-V & \mbox{si}   & 3\pi/2<\theta\leq 2\pi \\
\end{array}
\right.
\]
\begin{enumerate}
\item Represente la condición de contorno gráficamente como función de $\theta$
y diga cuál es el período fundamental de la misma.
\item Calcule $\Phi$ en el interior del cilindro. 
\end{enumerate}

\vspace{0.5cm}
{\bf Soluci\'on }{\em 4)}
\vspace{0.5cm}

El potencial electrostático $\Phi(r,\theta,z)$ satisface la ecuación de Laplace
en el interior de un cilindro infinitamente largo de radio $b$. Como la
condición de contorno no depende de la coordenada $z$, la solución tampoco
dependerá de ella. Por lo tanto,

\[
\Phi=\Phi(r,\theta),
\]

y la ecuación de Laplace se reduce a

\[
\frac{1}{r}\frac{\partial}{\partial r}
\left(r\frac{\partial\Phi}{\partial r}\right)
+
\frac{1}{r^2}
\frac{\partial^2\Phi}{\partial\theta^2}
=0,
\qquad 0\le r<b.
\]

La condición de borde es

\[
\Phi(b,\theta)=
\begin{cases}
 V,&0<\theta\le\pi/2,\\[1ex]
-V,&\pi/2<\theta\le\pi,\\[1ex]
 V,&\pi<\theta\le3\pi/2,\\[1ex]
-V,&3\pi/2<\theta\le2\pi.
\end{cases}
\]

\subsection*{a) Representación gráfica y período fundamental}

La condición de contorno corresponde a una onda cuadrada que alterna entre los valores $+V$ y $-V$ cada intervalo de longitud $\pi/2$ (ver Fig.~\ref{fig:boundary_condition}).

\begin{figure}[ht]
    \centering
    \includegraphics[width=0.5\textwidth]{fig-4-1.pdf}
    \caption{Condiciones de contorno en la superficie del cilindro visto desde arriba.}
    \label{fig:boundary_condition}
\end{figure}


Puede escribirse de forma compacta como

\[
\boxed{
\Phi(b,\theta)=
V\,\operatorname{sgn}(\sin2\theta).
}
\]

Además,

\[
\Phi(b,\theta+\pi)
=
\Phi(b,\theta),
\]

mientras que

\[
\Phi(b,\theta+\pi/2)
=
-\Phi(b,\theta).
\]

Por lo tanto, el período fundamental de la función es

\[
\boxed{T=\pi.}
\]

{\bf b) Potencial en el interior del cilindro}

Buscamos una solución por separación de variables,

\[
\Phi(r,\theta)=R(r)\Theta(\theta).
\]

Sustituyendo en la ecuación de Laplace,

\[
\frac{r^2R''+rR'}{R}
=
-\frac{\Theta''}{\Theta}
=n^2,
\]

donde $n$ es la constante de separación.

Las ecuaciones resultantes son

\[
\Theta''+n^2\Theta=0,
\]

y

\[
r^2R''+rR'-n^2R=0.
\]

Sus soluciones generales son

\[
\Theta(\theta)
=
A_n\cos(n\theta)
+
B_n\sin(n\theta),
\]

y

\[
R(r)
=
C_nr^n
+
D_nr^{-n}.
\]

Como el potencial debe permanecer finito en el origen, debe ser

\[
D_n=0.
\]

Además, el valor medio de la condición de contorno es nulo, por lo que no
aparece un término constante.

En consecuencia,

\[
\Phi(r,\theta)
=
\sum_{n=1}^{\infty}
\left(\frac{r}{b}\right)^n
\left[
A_n\cos(n\theta)
+
B_n\sin(n\theta)
\right].
\]

Para determinar los coeficientes utilizamos la serie de Fourier de la onda
cuadrada,

\[
\operatorname{sgn}(\sin x)
=
\frac{4}{\pi}
\sum_{k=0}^{\infty}
\frac{\sin[(2k+1)x]}{2k+1}.
\]

Tomando $x=2\theta$,

\[
\Phi(b,\theta)
=
\frac{4V}{\pi}
\sum_{k=0}^{\infty}
\frac{\sin[(4k+2)\theta]}{2k+1}.
\]

Comparando esta expresión con la solución general, se concluye que sólo
contribuyen los armónicos $n=4k+2$, obteniéndose finalmente

\[
\boxed{
\Phi(r,\theta)
=
\frac{4V}{\pi}
\sum_{k=0}^{\infty}
\frac{1}{2k+1}
\left(\frac{r}{b}\right)^{4k+2}
\sin[(4k+2)\theta],
\qquad 0\le r<b.
}
\]

Esta función satisface la ecuación de Laplace en el interior del cilindro y
reproduce la condición de contorno al tomar el límite $r\rightarrow b^{-}$.





 
\end{document}


% Sobras
\prob{} Calcule $\exp(A)$, donde $A$ es la matriz
\[
A = \left(\begin{array}{rrr}
-2 & 0 & 1 \\
-1 & -3 & 3 \\
-1 & -1 & 1
\end{array}\right)
\]
{\em Ayuda:} Transforme $A$ a forma de Jordan, exponencie y antitransforme.




